\chapter{Herramientas MCP}
\label{sec:mcp-tools}

\section{Protocolo}

El servidor MCP (Model Context Protocol) expone \textbf{55 tools}
organizadas en 10 dominios. Cada tool se define con la interfaz
\code{Tool<I, O>} en \filepath{platform/mcp/src/lib/types.ts}:

\begin{lstlisting}[language=TypeScript]
export interface Tool<I = unknown, O = unknown> {
  readonly name: string;
  readonly domain: string;
  readonly description: string;
  readonly stability: 'stable' | 'experimental' | 'deprecated';
  readonly version: string;
  readonly inputSchema: z.ZodType<I>;
  readonly outputSchema?: z.ZodType<O>;
  run(input: I, ctx: ToolContext): Promise<O>;
}
\end{lstlisting}

Cada tool es un archivo en \filepath{platform/mcp/src/tools/<dominio>/<nombre>.ts}
y se registra vía el \filepath{index.ts} del dominio. Los 10
dominios están listados en las secciones siguientes.

\section{Dominio \code{agentic} (9 tools)}

Flujo agentic slow-lane: negociación de ingesta, generación de
parser asistida, revisión humana, publicación de mappings
compartidos.

\begin{itemize}
    \item \code{agentic_review_list} — lista propuestas pendientes.
    \item \code{agentic_review_approve} — aprueba.
    \item \code{agentic_review_reject} — rechaza.
    \item \code{canonical_map} — normaliza eventos vía
    \code{FieldConfig}.
    \item \code{fetch_shared_mappings} — trae mappings del library.
    \item \code{field_config_schema} — schema actual de
    \code{FieldConfig}.
    \item \code{ingest_negotiate} — handshake de ingesta.
    \item \code{parser_generate} — infiere \code{FieldConfig}.
    \item \code{publish_mapping} — publica un mapping validado.
\end{itemize}

\section{Dominio \code{catalogs} (4 tools)}

\begin{itemize}
    \item \code{get_catalog} — catálogo de entity/relation types.
    \item \code{get_catalog_diagnostics} — checks de integridad.
    \item \code{get_catalog_with_status} — catálogo con banderas de
    salud.
    \item \code{list_datasets} — datasets disponibles.
\end{itemize}

\section{Dominio \code{data-quality} (7 tools)}

\begin{itemize}
    \item \code{invariant_checker} — corre invariantes del grafo.
    \item \code{invariant_check_hypothetical} — dry-run antes de
    ingestar.
    \item \code{mapping_regression_test} — prueba mapping contra
    fixture congelado.
    \item \code{schema_drift_detect} — compara batch contra esquema
    conocido.
    \item \code{list_dead_letters} — lista DLQ.
    \item \code{reingest_dead_letter} — reintenta evento.
    \item \code{purge_dead_letters} — borra entradas DLQ.
\end{itemize}

\section{Dominio \code{export} (3 tools)}

\begin{itemize}
    \item \code{export_iocs} — IOCs STIX-like.
    \item \code{export_ocsf} — OCSF v1.4.
    \item \code{export_subgraph} — subgrafo como GraphSON/JSONL.
\end{itemize}

\section{Dominio \code{graph-views} (4 tools)}

\begin{itemize}
    \item \code{get_entity_type_counts} — histograma de tipos.
    \item \code{get_graph_summary} — conteos globales.
    \item \code{get_subnet_analysis} — vista por subnet.
    \item \code{get_temporal_heatmap} — matriz hora $\times$ día.
\end{itemize}

\section{Dominio \code{hunting} (6 tools)}

\begin{itemize}
    \item \code{run_hunt} — ejecuta hipótesis DSL.
    \item \code{get_hunt_results} — resultados persistidos.
    \item \code{diff_hunts} — diferencia entre dos hunts.
    \item \code{export_hunt_results} — CSV/JSON.
    \item \code{parse_dsl} — valida DSL sin ejecutar.
    \item \code{save_hypothesis_as_sigma} — genera regla Sigma.
\end{itemize}

\section{Dominio \code{integrations} (3 tools)}

\begin{itemize}
    \item \code{enrich_ip} — GeoIP/WHOIS de una IP.
    \item \code{file_ticket} — abre ticket en el sistema externo.
    \item \code{publish_iocs_to_sentinel} — empuja IOCs a Sentinel.
\end{itemize}

\section{Dominio \code{navigation} (10 tools)}

\begin{itemize}
    \item \code{search_entities} — búsqueda por valor/tipo.
    \item \code{expand_node} — vecinos de un nodo.
    \item \code{expand_node_grouped} — vecinos agrupados por
    relación.
    \item \code{check_connection} — ping + auth check.
    \item \code{get_entity_types} — lista de tipos.
    \item \code{get_events_for_node} — eventos del nodo.
    \item \code{get_events_paginated} — eventos paginados.
    \item \code{get_node_details} — propiedades + eventos.
    \item \code{get_relation_schema} — matriz
    \code{(from, rel, to)}.
    \item \code{get_subgraph} — subgrafo inducido.
\end{itemize}

\section{Dominio \code{notes} (6 tools)}

\begin{itemize}
    \item \code{create_note} — crea nota sobre entidad/path/hunt.
    \item \code{get_notes} — lista.
    \item \code{tag_entity} — etiqueta.
    \item \code{untag_entity} — remueve etiqueta.
    \item \code{list_tagged} — entidades por etiqueta.
    \item \code{get_path_nodes} — nodos de un path concreto.
\end{itemize}

\section{Dominio \code{scoring} (3 tools)}

\begin{itemize}
    \item \code{enable_anomaly_scoring} — toggle por sesión.
    \item \code{compute_scores} — recomputa scoring sobre
    resultados existentes.
    \item \code{explain_score} — desglose de los 5 componentes para
    un path.
\end{itemize}

\section{Contrato común}

Todas las tools MCP comparten:

\begin{itemize}
    \item Autenticación: la API key del MCP se deriva de
    \code{GRAPHHUNTER_API_TOKEN}; el server MCP actúa como proxy
    tipado del API HTTP.
    \item Errores: cualquier error del API subyacente se propaga
    como \code{\{error: code, message, details\}} estructurado.
    \item Timeouts: 30s default para tools de lectura, 300s para
    tools de ingesta / hunt.
    \item Idempotencia: las tools marcadas como \code{stable} son
    seguras de reintentar con el mismo input.
\end{itemize}

\section{Tests contractuales}

El directorio \filepath{platform/mcp/test/} contiene tests que
validan cada tool contra un \filepath{platform/api} mock según el
schema del endpoint correspondiente. Un mismatch entre la definición
de la tool y el DTO del API rompe el build antes de release.
